\section{Polynomial-time exact decision for fixed species count}
\label{sec:fixed-n-independent}

Fix $n$ and define
\[
 K=\{v\in\mathbb R_{\geq0}^m:\Gamma v=0\},\qquad
 \Phi(v)=\Gamma\operatorname{diag}(v)Y^T,\qquad P=\Phi(K).
\]

\begin{lemma}[Relative-interior projection]
\label{lem:ri-projection-independent}
If $K$ contains a strictly positive vector, then
\[
 \{\Phi(v):v>0,\ \Gamma v=0\}=\operatorname{ri}P.
\]
\end{lemma}

\begin{proof}
A coordinate inequality $v_r\geq0$ is not identically tight on $K$ because a
strictly positive point exists. Hence
$\operatorname{ri}K=K\cap\mathbb R_{>0}^m$. For any nonempty convex set and
linear map $T$, one has $T(\operatorname{ri}C)=\operatorname{ri}T(C)$; this
follows by restricting $T$ to the affine hull and using the finite-dimensional
open-mapping theorem. Apply this with $T=\Phi$.
\end{proof}

\begin{lemma}[Circuit enumeration]
\label{lem:circuits-independent}
Every extreme ray of $K$ has support at most
$\operatorname{rank}\Gamma+1\leq n+1$. For fixed $n$, all extreme rays and
their exact rational generators are computable in polynomial time in the
reaction count and input bit length.
\end{lemma}

\begin{proof}
An extreme ray is a support-minimal nonzero nonnegative dependency among the
columns of $\Gamma$. If its support contained more than
$\operatorname{rank}\Gamma+1$ columns, the restricted kernel would have
dimension at least two. A nonproportional kernel perturbation can be scaled in
both directions until one coordinate vanishes while retaining
nonnegativity, decomposing the original vector into two nonproportional
vectors of $K$, a contradiction. Thus every support is a positive circuit of
size at most $n+1$.

Enumerate all column subsets of size at most $n+1$, compute their exact
nullspaces, and retain the one-dimensional dependencies with a same-sign
representative whose support is minimal. This includes singleton zero columns
and tolerates repeated columns. There are $O(m^{n+1})$ subsets. Circuit
coordinates can be chosen as signed minors of matrices of order at most $n$;
for fixed $n$, determinant bounds give polynomial encoding length.
\end{proof}

\begin{theorem}[Fixed-species complexity]
\label{thm:fixed-n-independent}
For every fixed number of species $n$, weak all-mobile stationary-Turing
admissibility of rationally encoded classical mass-action reaction networks is
decidable in polynomial bit complexity. The same is true for each fixed
designated mobile set.
\end{theorem}

\begin{proof}
First decide whether $K$ has a strictly positive point by rational linear
programming; by homogeneity this is equivalent to feasibility of
$\Gamma v=0$, $v_r\geq1$. If not, no positive equilibrium exists.

Map every circuit generator from Lemma~\ref{lem:circuits-independent} through
$\Phi$. Their images generate $P$ completely. The ambient dimension is
$n^2$, fixed. In fixed dimension, exact generator-to-relative-facet conversion
is polynomial: compute the span and lineality, quotient by the lineality,
enumerate candidate supporting hyperplanes from bounded-size subsets of
generators, and test them against all generators. Determinantal normal vectors
have polynomial bit length, and the fixed-dimensional upper bound makes the
number of candidates polynomial. Lifting back gives
\[
 E\operatorname{vec}(A)=0,\qquad \ell_j(A)>0
\]
for $A\in\operatorname{ri}P$. This description also covers lower-dimensional
and nonpointed images, $P=\{0\}$, and the case where the quotient has no
facets.

By Lemma~\ref{lem:ri-projection-independent}, admissibility is equivalent to
an existential formula in the $n^2$ entries of $A$ and the $n$ positive
entries of $h$: the relative-facet conditions above, the Routh--Hurwitz
conditions for $A\operatorname{diag}(h)$ restricted to
$S=\operatorname{im}\Gamma$, and one negative signed principal minor of $A$.
A rational basis of $S$ and left inverse have polynomial bit length by minor
bounds. For fixed $n$, the number of real variables and all determinant and
Hurwitz degrees are fixed, while the number and coefficient lengths of input
polynomials are polynomial. Fixed-variable algorithms for the existential
theory of the reals therefore run in polynomial bit complexity.

For a designated mobile set, add one real variable $\lambda$ and the relevant
principal-block determinant condition. The variable count remains fixed.
\end{proof}

This is an algorithmic theorem, not a finite motif or topology classification.
For $n=1,2,3$ it gives complete polynomial-time exact decision algorithms for
arbitrary reaction lists.
